% Extracto íntegro por residencia del propietario científico definitivo de masas.
\section{Équivalence entre la section persistante et sa représentation finie}

La définition généalogique de la particule et sa représentation opératorielle sont
deux présentations du même objet. Soit
\[
 \mathfrak r_K=(x_K,\gamma_K,B_K,s_K,n_K,U_K)
\]
un représentant fini à la profondeur \(K\), avec état, route, frontière,
orientation, compteur et représentation, et soient
\(r_{L,K}:\mathfrak R_L\to\mathfrak R_K\) les applications de raffinement. Une
section persistante est une famille
\[
 X=(\mathfrak r_K)_{K\ge K_0},
 \qquad r_{L,K}(\mathfrak r_L)=\mathfrak r_K
 \quad(L\ge K\ge K_0).
\]
Deux familles sont équivalentes si elles coïncident sur une queue cofinale.

\begin{theorem}[Équivalence de la section persistante et de la limite projective]
La catégorie des sections persistantes, définies comme des familles compatibles,
est canoniquement équivalente à
\[
 \varprojlim_K\mathfrak R_K/\!\sim_{\rm cola}.
\]
La projection à une profondeur finie est une évaluation, et non un foncteur inverse.
Un représentant \(\mathfrak r_K\) détermine une unique section persistante si
et seulement s'il possède un unique prolongement compatible modulo l'équivalence de queue ;
cette unicité est obtenue, par exemple, lorsque les applications de raffinement sont
bijectives sur la queue survivante à partir de \(K\).
\end{theorem}

\begin{proof}
Par définition, une section persistante est précisément un point de la limite
projective ; le quotient identifie les familles qui coïncident sur une queue cofinale.
L'évaluation est la projection canonique sur une composante. Sa fibre contient
tous les prolongements compatibles du représentant évalué, de sorte qu'elle
est unitaire exactement sous la condition de prolongement unique indiquée.
\end{proof}

Une particule HMT--MD est donc une classe de sections persistantes sans
obstruction terminale finie, munie d'une représentation, d'une orientation, d'une charge,
d'un spin et d'un caractère spectral compatibles avec le raffinement. La route est son
histoire observable, et non son identité réduite à une cellule. L'antiparticule
est la section conjuguée par l'involution dodécaphasique et la représentation
conjuguée. Une section finie obstruée est virtuelle ; un produit cocyclique de
sections est composé ; une représentation tressée persistante est anyonique.

\section{Théorème de clôture typée de la réalisation massique}

\begin{theorem}[Clôture par codomaine physique]
Supposons construits l'atlas \(81/56/13/324\), le secteur nul, l'échelle
action--horloge--masse, les lecteurs \(K_D,K_\Omega\), le descripteur physique et la
représentation interne de chaque état. Alors tout descripteur admissible et
réalisable, c'est-à-dire tout \(X\) avec
\(d_{\mathcal P}(X)\in\operatorname{im}d_{\mathcal R}\),
détermine, sans consulter sa valeur externe :
\begin{enumerate}
\item une fibre, une orbite ou une multisection \(\mathscr S(X)\) ;
\item la paire positive bidirectionnelle et sa PVM conjointe
      \((\widehat M_X^{\beta,D},\widehat M_X^{\beta,\Omega},
      E_X^{D,\Omega})\) ;
\item une décomposition en blocs irréductibles ;
\item une mesure spectrale pour chaque couplage physique ;
\item un scalaire dans un secteur comprimé de rang un ou un multiplet dans un
      secteur réductible ; dans un secteur ouvert, un opérateur effectif et, seulement s'il
      existe un prolongement méromorphe avec un zéro isolé admissible, un pôle
      complexe ;
\item un objet catégorique lorsque la réalisation est topologique et ne possède pas encore
      de gap hamiltonien ;
\item un critère d'inexistence lorsque le type demandé contredit
      la positivité, la jauge ou l'asymptoticité.
\end{enumerate}
Aucun de ces résultats ne requiert de sélectionner une route au moyen de la masse
qui sera ensuite confrontée.
\end{theorem}

\begin{proof}
Le produit fibré construit la fibre de manière prospective. Le caractère et les
deux lecteurs construisent la paire positive et sa PVM conjointe. Lorsque le type
physique déclare une fonctionnelle scalaire, le calcul fonctionnel produit l'opérateur
correspondant ; l'espérance conditionnelle le projette sur le commutant de la
symétrie. La décomposition isotypique produit des scalaires ou des blocs par le lemme
de Schur. Les projecteurs de préparation produisent des mesures spectrales. La
réduction de Feshbach produit l'opérateur effectif ; le prolongement méromorphe
et un zéro isolé admissible produisent un pôle, s'il existe. Les secteurs
topologiques conservent leurs données de fusion et de tressage jusqu'à ce qu'il existe une
réalisation hamiltonienne. La positivité et les contraintes de jauge
excluent les types incompatibles. Chaque étape dépend uniquement de données internes ou
du couplage déclaré, jamais de la colonne externe.
\end{proof}

\section{Conditions de cohérence et d'incompatibilité des réalisations spectrales, fermioniques et composées}

La construction est réfutée dans son domaine si l'un quelconque des
faits suivants se produit :

\begin{enumerate}
\item la modification des valeurs externes change la fibre \(\mathscr S(X)\) ;
\item l'attribution cesse d'être équivariante sous l'inversion cellulaire,
      la conjugaison ou la jauge ;
\item une cellule ponctuelle est sélectionnée dans une multisection non centrale sans
      couplage brisant la symétrie ;
\item la synthèse de \(K_D\) et \(K_\Omega\) privilégie une orientation ou ne
      retrouve pas la diagonale lorsque les deux lecteurs coïncident ;
\item une compression à plusieurs valeurs propres est remplacée par une moyenne sans
      projecteur physique ;
\item le quotient de couleur produit des masses différentes pour le rouge, le vert et le bleu ;
\item des masses de quarks sont comparées sans déclarer le même schéma et la même échelle ;
\item une phase de rang un est présentée comme mélange PMNS après avoir échoué à la
      confrontation angulaire de rang complet ;
\item la neutralité est utilisée à elle seule pour affirmer le caractère Majorana ;
\item la masse d'un composé est obtenue en additionnant les masses des constituants sans
      registre de liaison ;
\item la largeur est insérée comme coordonnée de la signature au lieu d'être extraite
      d'un mode ouvert ou d'un pôle ;
\item une valeur propre négative du Hessien est publiée comme particule tachyonique ;
\item la gravitation est rendue dépendante d'un graviton élémentaire introduit
      avant la connexion, la courbure et la torsion.
\end{enumerate}

\section{Confrontation métrologique et archive exhaustive des réalisations}

La comparaison externe n'est ouverte qu'après avoir figé le descripteur,
la fibre, l'opérateur, le projecteur, le schéma et l'échelle. Chaque ligne doit distinguer la masse
stable, la masse de pôle, l'estimateur de résonance et la largeur. L'atlas de \(324\)
fiches et les inventaires de \(471\) masses et de \(384\) largeurs constituent
l'exposé exhaustif qui suit. La construction conceptuelle et chaque fiche
renvoient réciproquement l'une à l'autre par leurs identifiants de route. Les fiches
conservent tous leurs champs et sont composées comme des unités cohérentes ; elles ne sont pas
remplacées par des tableaux historiques abrégés.

La comparaison ne consiste pas à élargir une liste de noms, mais à appliquer,
pour chaque type physique, l'application complète qui conduit de la généalogie de route
à son codomaine observable.
